\documentclass[a4paper]{article}

\input{patterns_L2_22}
%%%% Packages %%%%
\usepackage{amsfonts}
\usepackage{amsmath}
\usepackage{amssymb}
\usepackage{amsthm}
\usepackage{bbm}
\usepackage{booktabs}
\usepackage{calc}
\usepackage{caption}
\usepackage{centernot}
\usepackage{circuitikz}
\usepackage{dsfont}
\usepackage{esint}
\usepackage{esvect}
\usepackage{float}
\usepackage{fullpage}
\usepackage{graphicx}
\usepackage{hhline}
\usepackage{hyperref}
\usepackage{ifthen}
\usepackage{lipsum}
\usepackage{marvosym}
\usepackage{mathrsfs}
\usepackage{mathtools}
\usepackage{microtype}
\usepackage{multicol}
\usepackage{parskip}
\usepackage{siunitx}
\usepackage{svg}
\usepackage{tikz}
\usepackage{tikz-cd}
\usepackage{titlesec}
\usepackage{titling}
\usepackage{wrapfig}
\usepackage{xcolor}
\usepackage[inline]{enumitem}
\usepackage[makeroom]{cancel}
\usepackage[most]{tcolorbox}
\usetikzlibrary{angles}
\usetikzlibrary{arrows.meta}
\usetikzlibrary{bending}
\usetikzlibrary{calc}
\usetikzlibrary{patterns}
\usetikzlibrary{decorations.pathmorphing}
\usetikzlibrary{decorations.markings}

\setsvg{inkscapeexe="C:/Program Files/Inkscape/bin/inkscape.exe"}

\hypersetup{
	pdfborder={0 0 0}
}
\makeatletter
\renewcommand{\@makefnmark}{\hbox{\textsuperscript{\color{red}\normalfont\@thefnmark}}}
\makeatother
%%%% Packages %%%%
%%%% Custom Commands %%%%

\DeclareMathOperator{\Div}{div}

% Define a problem counter that resets with the section
\newcounter{problem}[section]

% Redefine the problem numbering to include section and problem numbers
\renewcommand{\theproblem}{\arabic{problem}}

% Define the Problem command - BG and ENG versions
\newcommand{\Problem}{\refstepcounter{problem}\fbox{\textbf{Problem \theproblem:}} }
\newcommand{\ProblemBG}{\refstepcounter{problem}\fbox{\textbf{Задача \theproblem:}} }

% Solution Definition
\newcommand{\Solution}[1]{\textbf{Solution\ifthenelse{\equal{#1}{}}{}{ #1}:}}

% This forces LaTeX to reserve vertical space as if every line had a fraction - use in \begin{aligned} .. \end{aligned}
\newcommand{\rowstrut}{\vphantom{\dfrac{1}{2}}}
%%%% Custom Commands %%%%
\tikzset{
	battery/.style={
		draw=none,
		decorate,
		decoration={
			markings,
			mark=at position 0.5 with {
				% Connectors
				\draw[line] (-\pgfdecoratedcompleteddistance,0) -- (-4,0);
				\draw[line] (4,0) -- (\pgfdecoratedcompleteddistance,0);
				% Battery
				\draw[line] (-4, -20) -- (-4, 20);
				\draw[line] (4, -10) -- (4, 10);
				% Polarity
				\node[above] at (-4, -15) {$+$};
				\node[below] at (4, -15) {$-$};
			}
		}
	}
}

\tikzset{
	switch/.style={
		draw=none,
		decorate,
		decoration={
			markings,
			mark=at position 0.5 with {
				% Connectors
				\draw[line] (-\pgfdecoratedcompleteddistance,0) -- (-15,0);
				\draw[line] (15,0) -- (\pgfdecoratedcompleteddistance,0);
				% Switch
				\draw[line] (-12.5,0) circle (2.5);
				\draw[line] (12.5, 0) circle (2.5);
				\draw[line] (-10,0) -- (10,7);
			}
		}
	}
}

\tikzset{
	spring/.style={
		decorate,
		decoration={
			coil,
			aspect=0.6,
			segment length=7pt,
			amplitude=6pt,
			pre length=5pt,
			post length=5pt
		},
		line width=1.5
	}
}

\tikzset{
	line/.style={
		line width=1.5pt,
	}
}

\tikzset{
	point/.style={
		circle,
		fill=black,
		inner sep=1.5pt
	}
}



%%% THIS IS A HELPER GRID FOR TIKZ PICTURES %%%

%\begin{center}
%	\begin{tikzpicture}
%		
%		% Grid
%		\draw[help lines, step=1] (-2,-2) grid (12,12);
%		
%		% Axes
%		\draw[->, line width=1] (-2,0) -- (12.5,0) node[right] {$x$};
%		\draw[->, line width=1] (0,-2) -- (0,12.5) node[above] {$y$};
%		
%		% Axis labels
%		\foreach \x in {-2,-1,1,2,...,12}
%		\node[below] at (\x,0) {\small $\x$};
%		
%		\foreach \y in {-2,-1,1,2,...,12}
%		\node[left] at (0,\y) {\small $\y$};
%		
%		% Your drawing
%		\draw[spring] (0,0) -- (10,10);
%		
%	\end{tikzpicture}
%\end{center}



%%% ROTATING ARROW HEADS %%%

% To rotate {Stealth} in place use the flex' command:
%\draw[
%line,
%-{Stealth[length=7pt,flex'=0.9]}] 
%(175,186.5) arc (90:270:12.5);



%%% EXAMPLE HATCHING FOR SURFACES %%%
%\begin{scope}
%	\clip (350,100) rectangle (425,275);
%	
%	\foreach \x in {175,183,...,425} {
%		\draw[line width=0.75pt]
%		(\x,100) -- ++(175,+175);
%	}
%\end{scope}

\usepackage[english]{babel}

\usepackage[
left=2.5cm,
right=2.5cm,
top=3cm,
bottom=2.5cm
]{geometry}

\setlength{\droptitle}{-7em}
\title{\Huge Physics Olympiad L2-22}
\author{BigChunguz24 - Last Edit 30/09/2026}
\date{}

\begin{document}
	\maketitle
	
	\Problem
	\textbf{\underline{Rolling Plank}}
	\\
	% A Guide To Physics Problems Part 1 - Problem 1.5
	A uniform thin rigid rod of mass $M$ is supported by two rotating rollers whose axes are separated by a fixed distance $a$. The rod is initially placed at rest asymmetrically at a distance $x_0$ away from the equilibrium position. The coefficient of friction between the bar and the rollers is $\mu$. Determine the displacement of the centre of mass of the rod $x(t)$ when:
	\begin{description}
		\item[$\bullet$] The rollers rotate in opposite direction as shown in the figure;
		\item[$\bullet$] The rollers rotate in reverse directions to the ones shown in the figure;
	\end{description}
	
	\begin{center}
	\begin{tikzpicture}[x=0.75pt,y=0.75pt,yscale=-1,xscale=1]
		% Rectangles
		\draw[line] (100,100) -- (400,100) -- (400,150) -- (100,150) -- cycle ;
		\draw[line] (150,225) -- (450,225) -- (450,275) -- (150,275) -- cycle ;
		
		% Wheels
		\draw[line] (325,175) circle (25);
		\draw[line] (175,175) circle (25);
		\draw[line] (325,300) circle (25);
		\draw[line] (175,300) circle (25);
		\node[point] at (325,175) {};
		\node[point] at (175,175) {};
		\node[point] at (325,300) {};
		\node[point] at (175,300) {};
		
		% Dashed Lines
		\draw [line, dashed] (250,100) -- (250,275) ;
		\draw [line, dashed] (300,225) -- (300,275) ;

		% Straight Lines
		\draw[line, {Stealth}-{Stealth}]    (250,250) -- (300,250) 
		node[midway, above] {$x_0$};
		
		\draw[line, {Stealth}-{Stealth}] (178,175) -- (322,175) 
		node[pos=0.45, above] {$a$};
		
		\draw[line, {Stealth}-{Stealth}] (178,300) -- (322,300) 
		node[midway, above] {$a$};
		
		% Arrows (curved)		
		\draw[
		line,
		-{Stealth[length=7pt,flex'=0.9]}] 
		(175,186.5) arc (90:270:12.5);
		
		\draw[
		line,
		{Stealth[length=7pt,flex'=0.9]}-] 
		(325,163.5) arc (-90:90:11.5);
		
		\draw[
		line,
		{Stealth[length=7pt,flex'=0.9]}-] 
		(175,311.5) arc (90:270:11.5);
		
		\draw[
		line,
		-{Stealth[length=7pt,flex'=0.9]}] 
		(325,288.5) arc (-90:90:11.5);
	\end{tikzpicture}
\end{center}
	
	\Problem
	\textbf{\underline{Interacting Dipoles}}
	\\
	% Irodov 3.39 + 3.40 + 3.41 + Addition
	Assume you are given a dipole with known charge $q$ and length $l$, or equivalently the dipole moment $\vv{p}$.
	\begin{description}[leftmargin=!, labelwidth=\widthof{$\bullet$}]
		\item[$\bullet$] Find an expression for the potential $\varphi$ of the dipole at point $S$ as shown on the figure below;
		\item[$\bullet$] Using this expression, find the polar components of the electric field strength vector;
		\item[$\bullet$] Find the points at which $\vv{E}$ is perpendicular to $\vv{p}$;
		\item[$\bullet$] The dipole is placed in an external uniform electric field whose strength is $\vv{E_0} \uparrow \uparrow \vv{p}$. In this case one of the equipotential surfaces enclosing the dipole forms a sphere. Determine the radius of this sphere;
		\item[$\bullet$] In figures $1)$ through $8)$ two identical dipoles are placed in different arrangements as shown, with their centres being a distance $d >> l$ from one another in each case. For each arrangement shown find approximated expressions for the energy of interaction $E_p$ and the force of interaction $\vv{F}$.
	\end{description}
	
	\begin{center}
	\begin{tikzpicture}[x=0.65pt,y=0.65pt,yscale=-0.99,xscale=0.99]
		% Nodes
		\node[point] at (25, 250) {};
		\node[point] at (25, 350) {};
		\node[point] at (136, 194) {};
		
		% Point S
		\draw[line, -{Stealth}] (25,300) -- (135,195) ;
		
		% Dipoles
		\draw [line, -{Stealth}]    (25,350) -- (25,253) ;
		\draw [line, -{Stealth}]    (225,210) -- (225,163) ;
		\draw [line, -{Stealth}]    (350,210) -- (350,163) ;
		\draw [line, -{Stealth}]    (225,285) -- (225,238) ;
		\draw [line, {Stealth}-]    (350,282) -- (350,235) ;
		\draw [line, -{Stealth}]    (225,360) -- (225,313) ;
		\draw [line, -{Stealth}]    (325,335) -- (372,335) ;
		\draw [line, -{Stealth}]    (225,435) -- (225,388) ;
		\draw [line, {Stealth}-]    (328,410) -- (375,410) ;
		\draw [line, -{Stealth}]    (475,400) -- (475,353) ;
		\draw [line, -{Stealth}]    (475,250) -- (475,203) ;
		\draw [line, -{Stealth}]    (550,400) -- (550,353) ;
		\draw [line, {Stealth}-]    (550,247) -- (550,200) ;
		\draw [line, -{Stealth}]    (625,400) -- (625,353) ;
		\draw [line, -{Stealth}]    (600,225) -- (647,225) ;
		\draw [line, {Stealth}-]    (679,225) -- (725,225) ;
		\draw [line, -{Stealth}]    (700,400) -- (700,353) ;
		
		% Dashed Lines
		\draw [line, dashed] (25,175) -- (25,425) ;
		\draw [line, dashed] (410,175) -- (410,425) ;
		\draw [line, dashed] (225,260) -- (350,260) ;
		\draw [line, dashed] (225,185) -- (350,185) ;
		\draw [line, dashed] (225,360) -- (350,360) ;
		\draw [line, dashed] (225,435) -- (350,435) ;
		\draw [line, dashed] (350,335) -- (350,360) ;
		\draw [line, dashed] (350,410) -- (350,435) ;
		\draw [line, dashed] (450,225) -- (450,375) ;
		\draw [line, dashed] (525,225) -- (525,375) ;
		\draw [line, dashed] (600,225) -- (600,375) ;
		\draw [line, dashed] (675,225) -- (675,375) ;
		\draw [line, dashed] (450,225) -- (475,225) ;
		\draw [line, dashed] (450,375) -- (475,375) ;
		\draw [line, dashed] (525,225) -- (550,225) ;
		\draw [line, dashed] (525,375) -- (550,375) ;
		\draw [line, dashed] (600,375) -- (625,375) ;
		\draw [line, dashed] (675,375) -- (700,375) ;
		\draw [line, dashed] (550,225) -- (550,375) ;
		\draw [line, dashed] (475,225) -- (475,375) ;
		\draw [line, dashed] (625,225) -- (625,375) ;
		\draw [line, dashed] (700,225) -- (700,375) ;
		
		% Text Node
		\draw (96,241.4) node [anchor=north west][inner sep=0.75pt]    {$r$};
		\draw (36,252.4) node [anchor=north west][inner sep=0.75pt]    {$\varphi $};
		\draw (135,190) node [anchor=south west][inner sep=0.75pt]    {$S$};
		\draw (186,182.4) node [anchor=north west][inner sep=0.75pt]    {$1)$};
		\draw (186,257.4) node [anchor=north west][inner sep=0.75pt]    {$2)$};
		\draw (186,332.4) node [anchor=north west][inner sep=0.75pt]    {$3)$};
		\draw (186,407.4) node [anchor=north west][inner sep=0.75pt]    {$4)$};
		\draw (466,426.4) node [anchor=north west][inner sep=0.75pt]    {$5)$};
		\draw (541,426.4) node [anchor=north west][inner sep=0.75pt]    {$6)$};
		\draw (616,426.4) node [anchor=north west][inner sep=0.75pt]    {$7)$};
		\draw (693,426.4) node [anchor=north west][inner sep=0.75pt]    {$8)$};
		\draw (281,167.4) node [anchor=north west][inner sep=0.75pt]    {$d$};
		\draw (281,241.4) node [anchor=north west][inner sep=0.75pt]    {$d$};
		\draw (281,341.4) node [anchor=north west][inner sep=0.75pt]    {$d$};
		\draw (281,416.4) node [anchor=north west][inner sep=0.75pt]    {$d$};
		\draw (436,291.4) node [anchor=north west][inner sep=0.75pt]    {$d$};
		\draw (511,291.4) node [anchor=north west][inner sep=0.75pt]    {$d$};
		\draw (586,292.4) node [anchor=north west][inner sep=0.75pt]    {$d$};
		\draw (661,292.4) node [anchor=north west][inner sep=0.75pt]    {$d$};
		
		%Shape: Arc [id:dp19626671896356962] 
		\draw  [draw opacity=0][line width=1.5]  (25.57,270.01) .. controls (33.55,270.15) and (40.77,273.42) .. (46.06,278.64) -- (25,300) -- cycle ; 
		\draw  [line width=1.5]  (25.57,270.01) .. controls (33.55,270.15) and (40.77,273.42) .. (46.06,278.64);  
	\end{tikzpicture}
\end{center}
	
	\Problem
	\textbf{\underline{Wave Superposition}}
	\\
	% Irodov 4.160
	Two plane waves $\psi_1$ and $\psi_2$ propagate without energy loss in a homogeneous elastic medium, one along the $x$-axis and the other along the $y$-axis: 
	\begin{description}
		\item[$\rightarrow$] $\psi_1(x, t) = a \cos(\omega t - k x)$;
		\item[$\rightarrow$] $\psi_2(x, t) = a \cos(\omega t - k y)$;
	\end{description}
	Find the wave motion pattern of the particles in the plane $xy$ if both waves are:
	\begin{description}
		\item[$\bullet$]  transverse and their oscillation directions coincide;
		\item[$\bullet$] longitudinal.
	\end{description}
	
	\Problem
	\textbf{\underline{Oscillating Capacitor}}
	\\
	% Задачи Московских Физических Олимпиад - 3.39
	Plate $A$ of a parallel-plate capacitor is fixed, while plate $B$ is attached to a wall with a spring and can move while remaining parallel to plate $A$. After closing the switch $K$, plate $B$ starts moving and eventually stops in a new equilibrium position. In this case, the initial equilibrium distance $d$ between the capacitor plates (when the spring is not stretched) decreases by $\eta = 10\%$. How much would the equilibrium distance between the plates change if the switch $K$ were closed for a short time? Assume that during this time, plate $B$ does not have time to move significantly.
	
	\begin{center}
	\begin{tikzpicture}[x=1pt,y=1pt,yscale=-1,xscale=1]
		% Left Capacitor Plate
		\draw[line]  (90,150) -- (100,150) -- (100,250) -- (90,250) -- cycle ;
		\begin{scope}
			\clip (90,150) rectangle (100,250);
			
			\foreach \x in {-100,-92,...,100} {
				\draw[line width=0.75pt]
				(\x,150) -- ++(100,100);
			}
		\end{scope}
		
		% Right Capacitor Plate
		\draw[line]  (150,150) -- (160,150) -- (160,250) -- (150,250) -- cycle ;
		\begin{scope}
			\clip (150,150) rectangle (160,250);
			
			\foreach \x in {50,58,...,154} {
				\draw[line width=0.75pt]
				(\x,150) -- ++(100,100);
			}
		\end{scope}
		
		% Hatching on the RHS
		\draw [line]    (325,150) -- (325,250) ;
		\begin{scope}
			\clip (325,150) rectangle (335,250);
			
			\foreach \x in {225,233,...,334} {
				\draw[line width=0.75pt]
				(\x,150) -- ++(100,100);
			}
		\end{scope}
		
		% Spring
		\draw[spring] (160, 200) -- (325, 200);
		
		% Connectors, Switch and Battery
		\draw[line] (95,330) -- (155,330) ;
		\draw[switch] (155,250) -- (155,330)
		node[midway, left] {$K$};
		\draw[battery] (95, 250) -- (95, 330);

		% Straight Lines and Dashed Lines
		\draw[line, {Stealth}-{Stealth}]    (100,125) -- (150,125) 
		node[midway, above] {$d$};
		\draw[line, dashed] (100,125) -- (100,150) ;
		\draw[line, dashed] (150,125) -- (150,150) ;
		
		% Text Node
		\draw (90, 150) node[above] {$A$};
		\draw (160, 150) node[above] {$B$};
	\end{tikzpicture}
\end{center}
		
	\Problem
	\textbf{\underline{Body-Pulley System}}
	\\
	% Новосибирски 2008 - 2.1.47
	Find the acceleration of the bodies in the system shown in the figure. The force $F$ is applied along the direction of the string to one of the bodies of mass $m$. The sections of the string on both sides of the light pulley, which is attached to a body of mass $M$, are parallel.
	
	\begin{center}
	\begin{tikzpicture}[x=1pt,y=1pt,yscale=-1,xscale=1]
		% Puley
		\draw[line] (186.835,150) circle (13.165);
		
		% Bodies
		\draw[line] (225,125) rectangle (275,175);
		\draw[line] (115,153.17) rectangle (135,173.17);
		\draw[line] (75,126.83) rectangle (95,146.83);

		% Straight Lines
		\draw[line] (186.83,150) -- (225,150) ;
		\draw[line] (186.83,163.17) -- (135,163.17) ;
		\draw[line] (186.83,136.83) -- (95,136.83) ;

		% Force
		\draw [line, -{Stealth}]    (75,136.8) -- (38,136.8) ;

		% Text Node
		\draw (35,122.4) node [anchor=north west][inner sep=0.75pt]    {$F$};

		\node at (85,137) {$m$};
		\node at (85,120) {$1$};
		
		\node at (125,163) {$m$};
		\node at (125,180) {$2$};
		
		\node at (250,150) {$M$};
		\node at (250,115) {$3$};
	\end{tikzpicture}
\end{center}
	
	\Problem
	\textbf{\underline{Battery Efficiency}}
	\\
	% 200 Puzzling Physics Problems - 155
	A battery consists of $N$ identical cells, each of e.m.f. $\varepsilon$. Is it true that the energy wasted when using the battery to charge a capacitor through a resistor can be reduced by charging it in $N$ stages? That is by connecting it first across a single cell, and then across two cells, and so on, rather than across the whole battery in a single step.

	
	\newpage
\vspace*{-5em}
\enlargethispage{\baselineskip}
\newcommand{\tg}{\operatorname{tg}}
\newcommand{\arctg}{\operatorname{arctg}}
\section*{\centering Appendix}
\section*{\centering Useful Mathematical Formulae and Physical Constants}
\emph{\underline{Important}}: Infinite series that are marked with $(*)$ are only valid under the condition that $|x| < 1$.

\begin{minipage}[t]{0.5\textwidth}
	\vspace{0pt}
	\begin{tcolorbox}[title=\textbf{Taylor Expansions}]
		\vspace{-10pt}
		\begin{align*}
			& \sin(x) = x - \frac{x^3}{6} + \frac{x^5}{120} + \mathcal{O}(x^7)
			&&
			\\[1pt]
			& \cos(x) = 1 - \frac{x^2}{2} + \frac{x^4}{24} - \frac{x^6}{720} + \mathcal{O}(x^8)
			&&
			\\[1pt]
			& \tg(x) = x + \frac{x^3}{3} + \frac{2x^5}{15} + \mathcal{O}(x^7)
			&&
			\\[1pt]
			& \arcsin(x) = x + \frac{x^3}{6} + \frac{3x^5}{40} + \mathcal{O}(x^7)
			&&
			\\[1pt]
			& \arctg(x) = x - \frac{x^3}{3} + \frac{x^5}{5} + \mathcal{O}(x^7)
			&&
			\\[1pt]
			& \exp(x) = 1 + x + \frac{x^2}{2} + \mathcal{O}(x^3)
			&&
			\\[1pt]
			& \ln(1+x) = x - \frac{1}{2}x^2 + \mathcal{O}(x^3)
			&& (*)
			\\[1pt]
			& (1-x)^{-1} = 1 + x + x^2 + \mathcal{O}(x^3)
			&& (*)
			\\[1pt]
			& (1+x)^{-1} = 1 - x + x^2 + \mathcal{O}(x^3)
			&& (*)
			\\[1pt]
			& (1+x)^{+\frac{1}{2}} = 1 + \frac{1}{2}x - \frac{1}{8}x^2 + \mathcal{O}(x^3)
			&& (*)
			\\[1pt]
			& (1+x)^{-\frac{1}{2}} = 1 - \frac{1}{2} x + \frac{3}{8} x^2 + \mathcal{O}(x^3)
			&& (*)
			\\[1pt]
			& (1+x)^{+\frac{3}{2}} = 1 + \frac{3}{2}x + \frac{3}{8}x^2 + \mathcal{O}(x^3)
			&& (*)
			\\[1pt]
			& (1+x)^{-\frac{3}{2}} = 1 - \frac{3}{2}x + \frac{15}{8}x^2 + \mathcal{O}(x^3)
			&& (*)
		\end{align*}
		
	\end{tcolorbox}
\end{minipage}
\begin{minipage}[t]{0.5\textwidth}
	\vspace{0pt}
	\begin{tcolorbox}[title=\textbf{Inverse Hyperbolic Functions}]
		\vspace{-10pt}
		\begin{align*}
			& \text{arsinh}(x) = \ln(x + \sqrt{x^2 + 1})
			\\[3.0375pt]
			& \text{arcosh}(x) = \ln(x + \sqrt{x^2 - 1})
			\\[3.0375pt]
			& \text{artgh}(x) = \frac{1}{2} \ln
			\left(
			\frac{1+x}{1-x}
			\right)
			\\[3.0375pt]
			& \text{arctgh}(x) = \frac{1}{2} \ln
			\left(
			\frac{x+1}{x-1}
			\right)
		\end{align*}
		
	\end{tcolorbox}
	\begin{tcolorbox}[title=\textbf{Indefinite Integrals}]
		\vspace{-10pt}
		\begin{align*}
			& \int (1-x^2)^{-\frac{3}{2}} dx = \frac{x}{\sqrt{1-x^2}} + C
			\\[3.0375pt]
			& \int (1+x^2)^{-\frac{3}{2}} dx = \frac{x}{\sqrt{1 + x^2}} + C
			\\[3.0375pt]
			& \int (1-x^2)^{-\frac{1}{2}} dx = \arcsin(x) + C
			\\[3.0375pt]
			& \int (1+x^2)^{-\frac{1}{2}} dx = \text{arsinh}(x) + C
			\\[3.0375pt]
			& \int (1-x^2)^{+\frac{1}{2}} dx = \frac{\arcsin(x) + x \sqrt{1-x^2}}{2} + C
			\\[3.0375pt]
			& \int (1+x^2)^{+\frac{1}{2}} dx = \frac{\text{arsinh}(x) + x \sqrt{1 + x^2}}{2} + C
		\end{align*}
		
	\end{tcolorbox}
\end{minipage}

\begin{minipage}{\textwidth}
	\begin{tcolorbox}[title=\textbf{Useful Physical Constants}]
		\vspace{-5pt}
		\begin{center}
			\begin{tabular}{|l|c|c|c|}
				\hline
				\textbf{Constant} & \textbf{Symbol} & \textbf{Value} & \textbf{Measuring Units} 
				\\
				\hline
				Light speed in vacuum & $c$ & $3{,}000 \times 10^8$ & $\mathrm{m/s}$
				\rule{0pt}{10pt}
				\\[1pt]
				\hline
				Rest mass of electron & $m_e$ & $9{,}109 \times 10^{-31}$ & $\mathrm{kg}$
				\rule{0pt}{10pt}
				\\[1pt]
				\hline
				Rest mass of proton & $m_p$ & $1{,}673 \times 10^{-27}$ & $\mathrm{kg}$
				\rule{0pt}{10pt}
				\\[1pt]
				\hline
				Rest mass of neutron & $m_n$ & $1{,}675 \times 10^{-27}$ & $\mathrm{kg}$
				\rule{0pt}{10pt}
				\\[1pt]
				\hline
				Rest mass of muon & $m_{\mu}$ & $1{,}884 \times 10^{-28}$ & $\mathrm{kg}$
				\rule{0pt}{10pt}
				\\[1pt]
				\hline
				Rest mass of tau-lepton & $m_{\tau}$ & $3{,}168 \times 10^{-27}$ & $\mathrm{kg}$
				\rule{0pt}{10pt}
				\\[1pt]
				\hline
				Elementary charge & $e$ & $1{,}602 \times 10^{-19}$ & $\mathrm{C}$
				\rule{0pt}{10pt}
				\\[1pt]
				\hline
				Planck's constant & $h$ & $6{,}626 \times 10^{-34}$ & $\mathrm{Js}$
				\rule{0pt}{10pt}
				\\[1pt]
				\hline
				Boltzmann's constant & $k$ & $1{,}381 \times 10^{-23}$ & $\mathrm{J/K}$
				\rule{0pt}{10pt}
				\\[1pt]
				\hline
				Gravitational constant & $\gamma$ & $6{,}673 \times 10^{-11}$ & $\mathrm{Nm^2/kg^2}$ 
				\rule{0pt}{10pt}
				\\[1pt]
				\hline
				Dielectric permittivity of vacuum & $\varepsilon_0$ & $8{,}854 \times 10^{-12}$ & $\mathrm{F/m}$
				\rule{0pt}{10pt}
				\\[1pt]
				\hline
				Avogadro's number & $N_A$ & $6{,}022 \times 10^{23}$ & $\mathrm{mol^{-1}}$
				\rule{0pt}{10pt}
				\\[1pt]
				\hline
				Rydberg's constant & $R_{\infty}$ & $1{,}097 \times 10^7$ & $\mathrm{m^{-1}}$
				\rule{0pt}{10pt}
				\\[1pt]
				\hline
				Magnetic permeability of vacuum & $\mu_0$ & $4\pi \times 10^{-7}$ & $\mathrm{H/m}$
				\rule{0pt}{10pt}
				\\[1pt]
				\hline
				Universal gas constant & $R$ & $8{,}314$ & $\mathrm{J / (mol \cdot K)}$
				\rule{0pt}{10pt}
				\\[1pt]
				\hline
			\end{tabular}
		\end{center}
	\end{tcolorbox}
\end{minipage}
	
	% A Guide To Physics Problems Part 1 - Problem 1.5
	% Irodov 3.39 + 3.40 + 3.41 + Addition
	% Irodov 4.160
	% Задачи Московских Физических Олимпиад - 3.39
	% Новосибирски 2008 - 2.1.47
	% 200 Puzzling Physics Problems - 155
\end{document}